\documentclass[11pt]{article}
\usepackage[a4paper,margin=1in]{geometry}
\usepackage{amsmath,amssymb,mathtools,bm}
\usepackage{enumitem,booktabs,array}
\usepackage{tcolorbox}
\usepackage{hyperref}
\usepackage[protrusion=true,expansion=false]{microtype}
\hypersetup{colorlinks=true,linkcolor=blue,urlcolor=blue,citecolor=blue}
\setlist[itemize]{topsep=2pt,itemsep=2pt}
\setlist[enumerate]{topsep=2pt,itemsep=2pt}
\newtcolorbox{keybox}{colback=blue!3!white,colframe=blue!40!black,boxrule=0.4pt,arc=1mm}
\newtcolorbox{alertbox}{colback=red!3!white,colframe=red!40!black,boxrule=0.4pt,arc=1mm}
\newtcolorbox{answerbox}{colback=green!3!white,colframe=green!40!black,boxrule=0.4pt,arc=1mm}
\newtcolorbox{drillbox}{colback=yellow!5!white,colframe=orange!60!black,boxrule=0.4pt,arc=1mm}
\newcommand{\EV}{\mathbb{E}}
\newcommand{\Var}{\mathrm{Var}}
\newcommand{\Cov}{\mathrm{Cov}}
\newcommand{\blank}[1]{\underline{\hspace{#1}}}

\title{E07-03: Agrawal \& Tambe (2016, RFS) \\
\large ``Private Equity and Workers' Career Paths: The Role of Technological Change'' --- Active Recall Workbook}
\author{Empirical Corporate Finance II --- Exam Prep}
\date{\today}
\begin{document}\maketitle

\begin{keybox}
\textbf{Paper in one sentence.} AT match LinkedIn-style resume data on millions of workers to PE-acquisition events and show that workers at PE-owned firms acquire more new-technology skills and accelerate their careers, with the effect concentrated in technology-exposed occupations and driven by PE-led IT investment --- evidence that PE ownership has long-run human-capital effects, not just financial restructuring.

\textbf{Placement.} Extends the PE-real-effects debate into labor economics. Connects with Davis-Haltiwanger-Jarmin-Lerner-Miranda (2014) on PE and jobs and with Bernstein-Sheen (2016) on operational channels.
\end{keybox}

\section*{Round 1: Complete Walk-Through}

\subsection*{1.1 Research question}
Most PE research measures firm-level outcomes (leverage, margins, closure). But PE could also affect workers' long-run human capital --- skills, mobility, wages. Does PE ownership accelerate technology adoption and therefore change worker career paths?

\subsection*{1.2 Data}
\begin{itemize}
\item Proprietary resume database covering $\sim$40M US workers.
\item CapitalIQ / Preqin PE-deal data, 2000--2014.
\item Worker-level skills coded by technology type (cloud, ERP, analytics).
\item Industry and occupation-level technology-exposure indices.
\end{itemize}

\subsection*{1.3 Empirical design}
\begin{itemize}
\item DiD: workers at firms that get PE-acquired vs.\ matched-firm workers.
\item Outcome: years-to-next-job, new-tech skill adoption, wage growth (where observed).
\item Interaction with occupation tech-exposure (IT-intensive vs.\ manual).
\end{itemize}

\subsection*{1.4 Headline results}
\begin{itemize}
\item Post-PE, workers add new-technology skills at $\sim$30--50\% higher rate than matched controls.
\item Career transitions accelerate: promotion probability rises, time-to-next-job falls.
\item Effect concentrated in tech-exposed occupations; minimal for low-tech roles.
\item Workers who stay at PE-owned firms see faster skill accumulation; leavers carry those skills to new employers.
\end{itemize}

\subsection*{1.5 Mechanism}
\begin{itemize}
\item PE owners invest in IT and modernize systems.
\item Workers learn new tools, raising human capital.
\item Labor-market spillovers: PE firms upgrade training for the broader industry.
\end{itemize}

\subsection*{1.6 Implications}
\begin{itemize}
\item PE effects are partly human-capital formation, not purely cost cutting.
\item Adds a positive-externality dimension to PE labor-market impact.
\item Speaks to the policy debate on PE and workers.
\end{itemize}

\subsection*{1.7 Caveats}
\begin{alertbox}
\begin{itemize}
\item Resume data is self-reported and may be biased toward professional workers.
\item Selection: PE acquires firms with growth prospects in tech-intensive industries.
\item Wage effects harder to measure without W-2 data.
\end{itemize}
\end{alertbox}

\section*{Round 2: Level 1 Fill-in}
\begin{enumerate}
\item Data: $\sim$\blank{1cm}M worker resumes.
\item Sample period: \blank{1cm}--\blank{1cm}.
\item Skill acquisition gain: \blank{1cm}--\blank{1cm}\%.
\item Effect concentrated in \blank{2cm}-exposed occupations.
\item Mechanism: PE-led \blank{1cm} investment.
\end{enumerate}

\section*{Round 3: Level 2 Prompts}
\begin{enumerate}
\item Describe the resume data and PE-deal matching.
\item Report findings on skill acquisition, career transitions, wages.
\item Explain the tech-exposure interaction.
\item Discuss mechanism and spillovers.
\item Place AT in the PE-labor-effect literature.
\end{enumerate}

\section*{Round 4: Minimal Scaffolding}
From a blank page: (i) question, (ii) data, (iii) findings, (iv) mechanism, (v) implications.

\section*{Round 5: Blank Page}
Essay: PE ownership and worker human-capital formation.

\vspace{6pt}
\begin{drillbox}
\textbf{Speed Drill.} (1) Data. (2) Treatment. (3) Skill effect. (4) Occupation pattern. (5) Channel.
\end{drillbox}

\section*{Quick Reference \& Answer Key}
\begin{answerbox}
\begin{itemize}
\item \textbf{Data.} $\sim$40M resumes + PE deals.
\item \textbf{Effect.} Tech-skill adoption $+30$--$50$\%.
\item \textbf{Concentration.} Tech-exposed occupations.
\item \textbf{Mechanism.} PE-led IT investment.
\end{itemize}
\end{answerbox}

\begin{keybox}
\textbf{30-second exam pitch.} Agrawal \& Tambe (2016) use a $\sim$40M-worker resume database combined with CapitalIQ / Preqin PE-deal data 2000--2014 to test whether PE ownership shapes worker human capital. In a DiD comparing workers at PE-acquired firms to matched controls, they find that post-acquisition, workers acquire new-technology skills at 30--50\% higher rates, see faster promotions and job transitions, and that skills persist when workers leave to new firms. Effects concentrate in technology-exposed occupations and reflect PE-led IT modernization. The paper adds a human-capital dimension to the PE-real-effects literature: PE is not purely financial restructuring or cost cutting; it also accelerates technology adoption and skill formation. Combined with Bernstein-Sheen (operational quality) and Davis et al.\ (employment dynamics), it supports a nuanced view that PE has multiple real-economy channels, some with positive spillovers to workers.
\end{keybox}

\end{document}
